
\documentclass[../../../physics-standard_questions_all.tex]{subfiles}

\begin{document}

真空中に定めた$x$軸に沿って，$x$軸の負の向きに大きさ$E$の平行一様電場が在る．
この電場中における正の点電荷（質量$m$，電荷$q$）の$x$軸に沿った運動を考える．
点電荷には，原点（$x=0$）において，点電荷に$x$軸の正の向きに初速$v_0$を与えたものとし，
クーロン力以外の力の影響は考えない．

\begin{enumerate}

    \item 点電荷が位置$x$に達したときの速度を$v$とする．
    そのときの速さ$|v|$を，$x$の関数として表せ．

    \item 位置$x$の電位$\phi(x)$を求めよ．
    ただし，$x=h$の位置を電位の基準とする．

\end{enumerate}

\end{document}


